% abstract.tex
The A=1 discrete causal lattice of Paper~I~\cite{menendez2026geometryfirst} is, by
construction, a birefringent medium: its bipartite hop geometry singles out the
optical axis $\hat{\mathbf{n}}_* = (1,1,-1)/\sqrt3$---the one cube body-diagonal
absent from the hop set---so propagation along the axis is inequivalent to
propagation across it. We work this structural fact to closed form in two sectors
that share that single axis. In the \emph{kinematic} sector, the single-tick spinor
propagator yields an exact ordinary/extraordinary speed split: the dispersion
depends only on the wavevector component perpendicular to $\hat{\mathbf{n}}_*$, with
effective-mass tensor $M_\text{eff} = \tfrac43 P_\perp$, a group-velocity split
$\Delta v = \tfrac23\sin\omega\,|\mathbf{k}|$ in the massive sector, and an
$\omega$-independent dispersion-curvature split ($4/3$ versus $0$) carried by the
photon and massive sectors alike. This is the paper's falsifiable dispersion
prediction: fixed in direction by $\hat{\mathbf{n}}_*$ and parameter-free up to the
lattice scale. In the \emph{gauge} sector, integrating the matter tokens out of the
A=1 dynamics induces an anisotropic Maxwell action whose permittivity and
inverse-permeability tensors, $\boldsymbol{\varepsilon}$ (eigenvalues $\{1,4,4\}$)
and $\boldsymbol{\mu}^{-1}$ (eigenvalues $\{4,4,16\}$)---both uniaxial about
$\hat{\mathbf{n}}_*$, in opposite sense---are read directly off the running engine's
hop holonomies. Because $\boldsymbol{\mu}^{-1} = \operatorname{adj}
(\boldsymbol{\varepsilon})$, the photon dispersion is a double transverse root and
the gauge-sector vacuum \emph{birefringence cancels exactly} (a null we confirm
numerically to $10^{-15}$), leaving a factor-$\sim2$ common-mode speed anisotropy
whose observability we leave open. The gauge-sector derivation and the adjugate
cancellation theorem are the companion Paper~VIII~\cite{menendez2026electricblock};
here they are confirmed independently from the engine. Every claim's scope, evidence,
and status is recorded in the audit table (Table~\ref{tab:audit}).
